\section{Comparison with alternatives}\label{sec:comparison}

Running models across more than one place is now a crowded field\footnote{Research as of 2026-11-21}, but
the efforts split into three camps: vendor hybrids that pair one personal device with the vendor's
cloud (Microsoft, Apple, Google), data-center stacks for operators of GPU fleets (NVIDIA Dynamo,
llm-d, GPUStack), and community tools that pool or serve models on hardware you own (exo, LocalAI,
LM Studio LM Link, Ollama, llama-swap, Petals). Ymir sits in the third camp but borrows the second
camp's scheduling ideas. Every claim below comes from the cited vendor or project page; Meta's
Llama Stack is left out because its current status could not be verified from a primary source.

{\small
\begin{xltabular}{\linewidth}{@{}L{3.3cm}L{3.2cm}YY@{}}
\caption{Who owns the hardware, pooling and exposure.}\\
\toprule
Product & Who owns the hardware & Pools several machines & Network exposure and auth \\
\midrule
\endfirsthead
\toprule
Product & Who owns the hardware & Pools several machines & Network exposure and auth \\
\midrule
\endhead
Ymir & You & Yes, any mix of your devices & No public endpoint; Weirdnet only, identity roster \\
Microsoft Hybrid Intelligence & You (one PC) plus Microsoft cloud & No; one PC and the cloud & Microsoft account, tenant identity \\
Microsoft Autopilot & Microsoft cloud & Not applicable (agent service) & Entra Agent ID inside the tenant \\
Apple Private Cloud Compute & You (device) plus Apple servers & No; device and Apple's servers & Apple-operated, attested nodes \\
Google Gemini Nano hybrid & You (phone) plus Google cloud & No; device and cloud fallback & Firebase project and Google cloud \\
NVIDIA Dynamo & Data-center operator & Yes, GPU fleets on Kubernetes & Operator's API gateway and auth \\
llm-d & Data-center operator & Yes, Kubernetes clusters & Kubernetes Gateway API \\
GPUStack & Operator (on-prem or cloud) & Yes, multi-cluster & REST API with auth, RBAC, tenants \\
exo & You & Yes, devices on one network & Same-network discovery by cluster name; local HTTP API \\
LocalAI federation & You or an operator & Yes, remotes over P2P invites & Invite links; admin key on the HTTP API \\
LM Studio LM Link & You & Remote access to your devices & Tailscale mesh, no public ports \\
Ollama & You & No; one machine & Local API with no inbound auth \\
llama-swap & You & No; one machine & Local OpenAI and Anthropic API \\
Petals & Volunteers or a private swarm & Yes, BitTorrent-style & Public swarm by default \\
\bottomrule
\end{xltabular}

\begin{xltabular}{\linewidth}{@{}L{3.1cm}L{2.9cm}L{2.7cm}L{2.6cm}Y@{}}
\caption{Failover, data locality, license and maturity, as of 2026-11-21.}\\
\toprule
Product & Mid-answer failover & Data locality & License & Maturity \\
\midrule
\endfirsthead
\toprule
Product & Mid-answer failover & Data locality & License & Maturity \\
\midrule
\endhead
Ymir & Yes, model and coordinator level & Only your linked devices & BSD-3-Clause & Pre-release; no tagged version \\
Microsoft Hybrid Intelligence & Not described & Local or Microsoft cloud, per task & Proprietary & Coming soon (Copilot); preview (Copilot Fusion) \\
Microsoft Autopilot & Not described & Microsoft cloud, in tenant & Proprietary & Rolling out to Frontier customers \\
Apple Private Cloud Compute & Not described & Device or Apple servers & Proprietary; code published for research & Shipping \\
Google Gemini Nano hybrid & Falls back to cloud & Device or Google cloud & Proprietary & Experimental (Firebase hybrid) \\
NVIDIA Dynamo & Yes, request migration & Operator's data center & Open source on GitHub & 1.0 since March 2026 \\
llm-d & Not described & Operator's cluster & Apache-2.0 & v1.0; CNCF Sandbox \\
GPUStack & Not described & Operator's clusters & Apache-2.0 & v2.3.1 (2026-11-11) \\
exo & Not described & Your devices & Apache-2.0 & v1.0.x releases \\
LocalAI federation & Failed requests move to the next member & Your instances & MIT & Active \\
LM Studio LM Link & Not described & Your devices & Not checked & Preview \\
Ollama & No & Your machine (cloud models optional) & MIT & Mature \\
llama-swap & No & Your machine & MIT & Mature \\
Petals & Not described & Shared swarm & MIT & Last commits in 2024 \\
\bottomrule
\end{xltabular}
}

\subsection{Microsoft Hybrid Intelligence and Autopilot}\label{sec:cmp-microsoft}

The closest large-company effort is Microsoft's Hybrid Intelligence for Windows, restated at
Ignite on 2026-11-17
(\url{https://blogs.windows.com/windowsexperience/2026/11/17/microsoft-ignite-2026-building-windows-for-the-next-era-of-work/}).
Copilot on Copilot+ PCs is to gain local context, local actions and local models, using models on
the PC and reaching to the cloud when a task needs it; Microsoft lists these as coming soon with
no date. GitHub Copilot Fusion already routes coding work between on-device models such as MAI
Code 1.1 Flash and cloud models, in experimental preview, and GitHub notes that routing decisions
are still made in the cloud
(\url{https://github.blog/changelog/2026-10-28-github-copilot-fusion-now-supports-selection-of-local-models/}).
Foundry Local and the Windows AI APIs supply the local models, including Phi Silica on Copilot+
PC NPUs.

Autopilot, shown in the Microsoft Mechanics video ``Introducing Autopilot, your personal and
proactive agent'' (2026-11-18), is a persistent agent that runs on dedicated Microsoft cloud
compute inside a company's tenant, with its own Entra Agent ID, and can run sessions on managed
devices; it began rolling out to Frontier customers on 2026-11-17. It is an agent product rather
than model infrastructure.

The difference with Ymir is where the cloud lives. Microsoft pairs one device with Microsoft's
data centers and decides per task which side runs it; nothing in these announcements pools
several of a user's own PCs into one shared model service. Ymir has no vendor side at all: the
cloud is the set of machines you own, linked to your identity, and failover goes to another of
your machines rather than to a provider. Microsoft's approach is stronger where the task needs
frontier-scale models, enterprise governance or zero setup; Ymir's is stronger where data must
never leave your hardware and you want several machines to act as one.

\subsection{Product notes}\label{sec:cmp-notes}

\begin{rowlist}
  \item[Apple Private Cloud Compute] Extends on-device Apple Intelligence to Apple-operated servers designed to be stateless and publicly verifiable, with binaries and a transparency log open to security researchers. Private, but on Apple's hardware, not yours.
  \item[Google Gemini Nano with Firebase AI Logic] Runs Gemini Nano on Android through AICore and falls back to a cloud Gemini model automatically. Experimental; on-device requests are single-turn, foreground only and limited to 4,000 tokens.
  \item[NVIDIA Dynamo] Data-center inference framework over vLLM, SGLang and TensorRT-LLM with disaggregated prefill and decode, KV-aware routing and autoscaling. Its request migration keeps tokens already generated and continues on a healthy worker, the closest analogue to Ymir's mid-answer resume.
  \item[llm-d] Kubernetes-native stack on vLLM and the Gateway API Inference Extension: prefix-cache-aware routing, prefill/decode disaggregation, KV offloading and autoscaling. CNCF Sandbox since March 2026.
  \item[GPUStack] Cluster manager that deploys vLLM, SGLang and TensorRT-LLM across on-prem, Kubernetes and cloud GPUs, with an OpenAI-compatible API, multi-tenancy, RBAC, monitoring and failure recovery.
  \item[exo] Connects your devices into one cluster with automatic discovery on the same network, then shards a model across them with tensor or pipeline parallelism, including RDMA over Thunderbolt 5 on Macs. It focuses on running models too big for one device; Ymir treats splitting as opt-in and routes whole models by default.
  \item[LocalAI federation] Serves one model name from several LocalAI instances joined by P2P invite links, routing least-busy or round-robin and moving a request to the next member when one fails. Also offers model-split workers over P2P.
  \item[LM Studio LM Link] Lets LM Studio on your devices use each other's models over a Tailscale mesh with no public ports. In preview. It is remote access to a chosen device rather than a pool with load balancing.
  \item[Ollama] The common single-machine server. Its local API has no inbound authentication, so exposing it means adding a proxy; optional Ollama cloud models need sign-in.
  \item[llama-swap] One Go binary that hot-swaps local model servers behind OpenAI and Anthropic APIs on one machine. Ymir's per-node supervisor does the same job and adds cross-machine routing.
  \item[Petals] Splits large models across a volunteer swarm. The project appears dormant, with its last repository activity in 2024.
\end{rowlist}

\subsection{Where Ymir differs, where others are stronger}\label{sec:cmp-verdict}

\begin{xltabular}{\linewidth}{@{}YY@{}}
\toprule
Ymir is different in & Others are stronger in \\
\midrule
\endhead
No public endpoint, proxy or API key; admission is your device roster & Maturity: Ollama, llama-swap, LocalAI and GPUStack ship tagged releases with large user bases \\
Mid-answer resume at two levels, for personal hardware & Raw scale and throughput: Dynamo and llm-d target GPU fleets with disaggregation and autoscaling \\
Any mix of your devices: NVIDIA, AMD, Apple, CPU-only, laptops and phones & Splitting one huge model: exo's topology-aware sharding over RDMA is further along than Ymir's opt-in llama.cpp RPC groups \\
Conversation context synced and encrypted across your own nodes & Frontier models and zero setup: Microsoft, Apple and Google hybrids \\
One shared client core for desktop and phone (Bragi) & Multi-tenant governance, RBAC and audit, which Ymir excludes by design \\
\bottomrule
\end{xltabular}
